% !PRINTED 637
\weberpage{637}
These formulae hold for every $\nu$, if we further stipulate with regard
to the $F_\nu$ that $F_h = F_k$ shall hold whenever $h \equiv k
\pmod{n-1}$. It then follows from 4. that the functions $F_0,
F_1 \ldots F_{n-2}$ remain unchanged by the substitution $s$, and
undergo a cyclic permutation by $t^{-1}$.

If we wish here to reduce the exponents to their least positive residues
with respect to the modulus $n$, we set

\weq{15}{g^\nu = n q_\nu + r_\nu, \quad 0 < r_\nu < n, \quad r_0 = 1,}

and the numbers $r_\nu$ coincide, in some order, with the numbers
$1, 2 \ldots n-1$, so that $r_0$ is always $= 1$.

Then by (5) we have

\weq{16}{f_{\nu-1} = (\varepsilon^{r_\nu}, x)(\varepsilon^{r_{\nu-1}}, x)^{-g},}

and we further set

\weq{17}{\Phi_\nu = F_\nu f_\nu^{q_{n-2}} f_{\nu+1}^{q_{n-3}} \cdots f_{\nu+n-3},}

and obtain from (14)

\weq{18}{(\varepsilon^{r_\nu}, x)^n = \Phi_\nu^n f_\nu^{r_{n-2}} f_{\nu+1}^{r_{n-3}} \cdots f_{\nu+n-2}^{r_0}.}

The exponents $r_\nu$ remain unchanged when $g$ is replaced by $g - l n$,
and are therefore independent of the condition (10) and \gesperrt{can be derived
from an arbitrary primitive root $g$}.

The characteristic properties of the functions $f_\nu, F_\nu, \Phi_\nu$
occurring in these formulae, whose totality we shall denote by $\omega_\nu$,
are, to repeat them once more, the following:

$\alpha)$ \gesperrt{By the cyclic permutation $s = (h, h+1)$ among the
indices of $x$ the functions $\omega_\nu$ are not changed.}

$\beta)$ \gesperrt{By the linear permutation $t^{-1} = (h, g^{-1} h)$ among the
indices of the $x$ the cyclic permutation $(0, 1, 2 \ldots n-2)$
is produced among the indices of the $\omega$.}

$\gamma)$ \gesperrt{By the substitution $\sigma = (\varepsilon, \varepsilon^g)$
the same cyclic permutation $(0, 1, 2 \ldots n-2)$
is effected among the indices of $\omega$.}

$\delta)$ \gesperrt{The cyclic functions of $\omega_\nu$ are metacyclic
functions of the $x$ and free of $\varepsilon$.}
